% !TEX root = ../eigenfaces_algorithm.tex
%=====================================================
\section{背景：我们想让机器“认人”，但它只会算数}

\subsection{问题是什么}

我们手里有 \textbf{Olivetti Faces} 人脸库：40 个人，每人 10 张灰度照片，
每张照片是 $64\times 64$ 个像素，每个像素是一个 $0\sim 1$ 之间的亮度数字。
任务有两个：

\begin{enumerate}
  \item \textbf{重建}：用很少的“信息”把一张人脸大致拼回来；
  \item \textbf{识别}：拿来一张没见过的人脸，判断它是 40 个人里的哪一位。
\end{enumerate}

\subsection{第一步：把照片变成向量}

计算机眼里没有“脸”，只有数字。把 $64\times 64$ 的像素网格按行拉直，
就得到一个 $4096$ 维的列向量：
\[
\mathbf{a}=\begin{pmatrix}a_1\\ a_2\\ \vdots\\ a_{4096}\end{pmatrix}\in\R^{4096}.
\]
于是“一个人”在这个数据库里的身份，变成了 $\R^{4096}$ 空间里的一个点。
40 个人 $\times$ 10 张照片，一共 400 个点。

\begin{intuition}
把每个人脸想成 $\R^{4096}$ 这个超高维空间里的一个“萤火虫”。
40 个人就是 40 簇萤火虫，每簇 10 只。识别问题就是：飞来一只新萤火虫，
看它离哪一簇最近。整个特征脸算法，本质上就是在想办法
“把灯光调暗，只留下萤火虫，让簇与簇分得更开”。
\end{intuition}

\subsection{为什么不能直接比较像素？}

最天真的做法：新照片 $\mathbf{a}$ 和库里的照片 $\mathbf{b}$ 算欧氏距离
$\norm{\mathbf{a}-\mathbf{b}}$，最近的就是答案。问题是：

\begin{itemize}
  \item 同一个人，光照从左边换成右边，两张照片的像素距离可能比\textbf{不同两个人}的距离还大；
  \item 4096 个维度里，绝大多数像素（背景、额头、脸颊）几乎是“死”的——不携带身份信息；
  \item 距离被这些无信息维度平均稀释，噪声淹没信号。
\end{itemize}

\begin{intuition}
辨认朋友靠的不是 4096 个像素点，而是几个关键特征：脸型轮廓、五官位置、发际线。
也就是说，\textbf{虽然照片活在 4096 维里，但“人的差异”只活在很少的几个维度里}。
特征脸算法的全部目标，就是自动找出这几个维度。
\end{intuition}
